Fix $n\in\mathbb{N}$ and distinct points $x_1,\dots,x_n\in\mathbb{R}^2$.

\begin{itemize}[leftmargin=*]
\item Convention: a \emph{circle} with center $x_i$ means the set $\{y\in\mathbb{R}^2: \|y-x_i\|=\rho\}$ for some radius $\rho>0$ (not the closed disk).
\item Hypothesis (``no centered circle contains three other points''): for every $i\in\{1,\dots,n\}$ and every radius $\rho>0$, the circle centered at $x_i$ of radius $\rho$ contains at most two points of $\{x_1,\dots,x_n\}\setminus\{x_i\}$. Equivalently,
\[
\forall i\ \forall \rho>0:\quad \bigl|\{j\neq i: \|x_i-x_j\|=\rho\}\bigr|\le 2.
\]
\item Let the \emph{distance set}
\[
D(x_1,\dots,x_n):=\{\|x_i-x_j\|: 1\le i<j\le n\},
\]
and write $|D|$ for its cardinality.
\end{itemize}

Question (as written): Does there exist an absolute constant $c>0$ and an $n_0\in\mathbb{N}$ such that for all $n\ge n_0$ and all point sets satisfying the hypothesis,
\[
|D(x_1,\dots,x_n)|\ \ge\ (1+c)\frac{n}{2}\ ?
\]
